\documentclass[10pt,a4paper]{amsart}
\usepackage[margin=19mm]{geometry}
\usepackage{amsmath,amssymb,mathtools}
\usepackage[scr=rsfs,cal=euler]{mathalfa}
\usepackage{xfrac}
\usepackage[unicode,hidelinks,bookmarks=false]{hyperref}
\newcommand{\ie}{i.e.\ }
\newcommand{\cf}{cf.\ }
\newcommand{\resp}{respectively\ }
\newcommand{\Subsection}{\subsection}
\providecommand{\nobreakdash}{\nobreak\hbox{-}}
\setlength{\emergencystretch}{2em}
\multlinegap0pt
% ---------------------------------------------------------------------------
% Editorial AMS wrapper prelude. The theorem environments and the mathematical macros
% below are transcribed verbatim from the paper's own preamble (cached-originals/source0.tex
% lines 51--82); the AMS amsart class, the named format packages of the pinned TeX Live
% runtime and this editorial formatting are not part of the author's text. No mathematics
% is changed.
% ---------------------------------------------------------------------------
\usepackage{bm}
\providecommand{\newblock}{\hskip .11em plus.33em minus.07em}

\newtheorem{definition}{Definition}[section]
\newtheorem{theorem}[definition]{Theorem}
\newtheorem{lemma}[definition]{Lemma}
\newtheorem{remark}[definition]{Remark}
\newtheorem{proposition}[definition]{Proposition}
\newtheorem{claim}[definition]{Claim}
\newtheorem{corollary}[definition]{Corollary}
\newcommand\R{\mathbb{R}}
\newcommand\Z{\mathbb{Z}}
\newcommand\B{\mathbb{B}}
\newcommand\C{\mathbb{C}}
\newcommand\N{\mathbb{N}}
\newcommand\T{\mathbb{T}}
\newcommand\mscH{\mathcal{H}}
\newcommand\mcaF{\mathcal{F}}
\newcommand\mcaI{\mathcal{I}}
\newcommand\mcaJ{\mathcal{J}}
\newcommand\mcaA{\mathcal{A}}
\newcommand\mcaB{\mathcal{B}}
\newcommand\mcaC{\mathcal{C}}
\newcommand\mcaD{\mathcal{D}}
\newcommand\mcaK{\mathcal{K}}
\newcommand\mcaH{\mathscr{H}}
\newcommand\mscP{\mathscr{P}}
\newcommand\mcaM{\mathcal{M}}
\newcommand\mcaR{\mathcal{R}}
\newcommand\mcaN{\mathcal{N}}
\newcommand\mcaP{\mathcal{P}}
\newcommand\mcaE{\mathcal{E}}
\numberwithin{equation}{section}

\begin{document}
\title[Decay estimates for discrete bi-Laplace operators with potentials]{Decay estimates for discrete bi-Laplace operators with potentials on the lattice $\Z$}
\author{Sisi Huang and Xiaohua Yao}
\address{Sisi Huang, Department of Mathematics, Central China Normal University, Wuhan, 430079, P.R. China}
\email{hss@mails.ccnu.edu.cn}
\address{Xiaohua Yao, Department of Mathematics, Key Laboratory of Nonlinear Analysis and Applications (Ministry of Education), and Hubei Key Laboratory of Mathematical Sciences, Central China Normal University, Wuhan, 430079, P.R. China}
\email{yaoxiaohua@ccnu.edu.cn}
\subjclass[2020]{Primary 34L40, 39A70, 81U24, 81U30}
\thanks{The work is partially supported by NSFC (No.12671120, No.12531005) and the Fundamental Research Funds for the Central Universities.}
\date{Source: arXiv:2506.23119v2, version of 2026; distributed under CC BY 4.0}
\keywords{Decay estimates, Discrete bi-Laplace operators, Classification of resonances, Asymptotic expansion, Beam equation}
\begin{abstract}
It is known that the discrete Laplace operator $\Delta$ on the lattice $\mathbb{Z}$ satisfies the following sharp time decay estimate:
$$\big\|e^{it\Delta}\big\|_{\ell^1\rightarrow\ell^{\infty}}\lesssim|t|^{-\frac{1}{3}},\quad t\neq0,$$
which is slower than the usual $ O(|t|^{-\frac{1}{2}})$ decay in the continuous case on $\mathbb{R}$. However, this paper shows that the discrete bi-Laplacian $\Delta^2$ on $\mathbb{Z}$ actually exhibits the same sharp decay estimate $|t|^{-\frac{1}{4}}$ as its continuous counterpart.

In view of the free decay estimate, we further investigate the discrete bi-Schr\"{o}dinger operators of the form $H=\Delta^2+V$ on the lattice space $\ell^2(\mathbb{Z})$, where $V$ is a class of real-valued decaying potentials on $\Z$. %satisfying $|V(n)|\lesssim \left<n\right>^{-\beta}$ for some $\beta>0$, with $\left<n\right>=(1+|n|^2)^{\frac{1}{2}}$.
%Our main results are threefold. First, we establish the limiting absorption principle for $H$. Second, we present a complete classification of resonance types of $H$ at the thresholds $0$ and $16$ and derive the corresponding asymptotic expansions of resolvent of $H$.
%Finally, under suitable decay conditions on $V$, we prove the following $\ell^1-\ell^{\infty}$ dispersive estimates for all these resonances types:
First, we establish the limiting absorption principle for $H$, and then derive the full asymptotic expansions of the resolvent of $H$ near the thresholds $0$ and $16$, including resonance cases. In particular, we provide a complete characterization of the different resonance types in $\ell^2$-weighted spaces. 

Based on these results above, we establish the following sharp $\ell^1-\ell^{\infty}$ decay estimates for all different resonances types of $H$ under suitable decay conditions on $V$:
$$\big\|e^{-itH}P_{ac}(H)\big\|_{\ell^1\rightarrow\ell^{\infty}}\lesssim|t|^{-\frac{1}{4}},\quad t\neq0,$$
where $P_{ac}(H)$ denotes the spectral projection onto the absolutely continuous spectrum space of $H$. 
Additionally, the  decay estimates for the evolution flow of discrete beam equation are also derived:
$$\|{\cos}(t\sqrt H)P_{ac}(H)\|_{\ell^1\rightarrow\ell^{\infty}}+\Big\|\frac{{\sin}(t\sqrt H)}{t\sqrt H}P_{ac}(H)\Big\|_{\ell^1\rightarrow\ell^{\infty}}\lesssim|t|^{-\frac{1}{3}},\quad t\neq0.$$
\end{abstract}

\maketitle
\noindent\emph{Source, version and licence.} \emph{Decay estimates for discrete bi-Laplace operators with potentials on the lattice $\Z$}, by Sisi Huang and Xiaohua Yao. The source of this editable unit is the e-print of \textbf{version v2} of arXiv:2506.23119, https://arxiv.org/abs/2506.23119v2, whose material is distributed under the Creative Commons Attribution 4.0 International licence, http://creativecommons.org/licenses/by/4.0/, as recorded in the arXiv licence metadata for this paper and shown on that abstract page. The whole of the body below is version v2, the newest version of the paper. Full rights notice: licenses/arxiv-2506.23119-CC-BY-4.0.txt.\par\smallskip
\noindent\textit{Editorial scope.} Counted unit: original Theorem 1.8 of the article, the statement carried by the source environment \texttt{theorem} with source label \texttt{Asymptotic expansion theorem}, cached-originals/source0.tex lines 396--443, printed as \textbf{Theorem 1.8} exactly as in the article. It is delivered together with the author's own complete proof as the single primary proof body, source0.tex lines 1852--2081 (230 lines): the \texttt{proof} environment that opens with ``Proof of Theorem \ref{Asymptotic expansion theorem} and part (i), runs through parts (ii)--(vi), and closes with ``This completes the proof.'' after part (vi), so all six expansions of the statement are proved. No proof marker is added and no mathematics is written, repaired or re-derived; the author's notation and the printed slips of the source are kept literal. Necessary complete context is retained uncounted, in source order: source0.tex lines 88--89 and 96--110, the author's own abstract (its hidden \texttt{comment} block at lines 90--95 is omitted); lines 126--128, the discrete Laplacian that the free-resolvent paragraph refers to; lines 351--369, Definitions 1.4 and 1.5, which fix $Q$, $v_k$, $\widetilde{Q}$, $\widetilde{v}_k$ and the projections $S_j$, $\widetilde{S}_j$ used throughout the counted proof; lines 383--394, Theorem 1.7, the resonance-type characterisation quoted by Remark 5.3; lines 471--505, the original section heading ``Limiting absorption principle'' with its subsection ``Free resolvent'' and the free-resolvent setup up to and including the definition of the weighted operator spaces $\B(s,s')$; lines 1597--1791, Lemma 5.1 (the Puiseux expansion of $R_0^{\pm}(\mu^4)$ at both thresholds) with its complete author proof; and lines 1793--1851, Lemma 5.2, Remark 5.3 and Lemma 5.4 (the Jensen--Nenciu reduction lemma used in Step 1 of every part of the counted proof), up to and including the sentence that announces the proof of Theorem 1.8. The complete author proof of Lemma 5.2, source0.tex lines 2084--2218, is not retained: only its statement is delivered, and every statement that the counted proof invokes is present in full. The article's own numbering is reproduced by explicit counter settings: the shared theorem-like counter declared by \texttt{\textbackslash newtheorem\{definition\}\{Definition\}[section]} and shared with theorem, lemma, remark, proposition, claim and corollary, and the section-scoped equation counter, are set before each retained passage, so the counted statement prints as Theorem 1.8, Theorem 1.7 as 1.7, Lemma 5.1 as 5.1, Lemma 5.2 as 5.2, Remark 5.3 as 5.3, Lemma 5.4 as 5.4, and the retained displays keep their original tags: (1.20)--(1.26) for the six expansions and the remainder estimate of Theorem 1.8, (2.1)--(2.3) for the free resolvent, and (5.1) and (5.3) for the two Puiseux expansions. Every \texttt{\textbackslash ref}, \texttt{\textbackslash eqref} and \texttt{\textbackslash cite} inside every retained passage resolves inside this delivered file; the only two citations are \texttt{\textbackslash cite[Lemma 2.1]\{JN01\}} and \texttt{\textbackslash cite[Lemma 3.1]\{KKK06\}}, delivered as the author's own bibliography entries transcribed verbatim from the source bibliography with the author's own printed numbers [32] and [34]. The e-print contains no \texttt{includegraphics} command anywhere and ships no artwork file: its single figure environment, source0.tex lines 525--592 with three \texttt{tikzpicture} drawings, lies outside every retained passage, so no figure environment, no drawing code and no graphics-inclusion command occurs anywhere in this delivered file and no artwork is redistributed. The source's own class is AMS \texttt{amsart} (\texttt{\textbackslash documentclass[reqno,11pt]\{amsart\}}); the e-print ships no class or style file, so no publisher class is shipped, used or redistributed, and the wrapper is the same AMS \texttt{amsart} class with the paper's own macros and theorem declarations transcribed from its preamble. Every declared body of this unit is an exact, unmodified span of source0.tex: no body is a bounded passage comparison, \texttt{omit} was not used anywhere, and consequently no fidelity receipt is asserted in delivery-evidence.json. The complete concatenated original source, with the paper's own bibliography, is bundled as sources/180682/source0.tex. Omitted from this editable unit and therefore absent from it, apart from the two bibliography entries named above, are source0.tex lines 1--87; 90--95; 111--125; 129--350; 370--382; 395--395; 444--470; 501--1596; 1792--1792; 2082--2937; among them the article's own preamble and title block, Definition 1.1, Theorem 1.2 (the main decay theorem), Remarks 1.3 and 1.6, the whole of Sections 3 and 4, Section 6 and Section 7, the appendix on commutator estimates and Mourre theory, the remaining 57 bibliography entries, and all tables and figures. The AMS \texttt{amsart} wrapper, the named format packages of the pinned TeX Live runtime and this editorial source, licence and scope paragraph are editorial formatting only.\par\smallskip
\setcounter{section}{1}
\setcounter{equation}{1}
\setcounter{definition}{0}
\begin{equation}\label{definition of Laplacian}
(\Delta_{\Z^d}\phi)(n):=\sum_{j=1}^d(\phi(n+e_j)+\phi(n-e_j)-2\phi(n)), \qquad\forall\ \phi\in\ell^2(\Z^d),
\end{equation}

\setcounter{section}{1}
\setcounter{equation}{19}
\setcounter{definition}{3}
\begin{definition}\label{definition of Sj}
Let $Q=I-P$~{\rm(}where $I$ is the identity operator{\rm)} and $v_k(n)=n^kv(n),k\in\N$. For any $j\in\{0,1,2,3\}$, define $S_j$ as the orthogonal projection onto the space $S_j\ell^2(\Z)$ with
\begin{itemize}
\item[{\rm(i)}]$S_0\ell^2(\Z):=\big\{f\in\ell^2(\Z):\left<f,v_{k}\right>=0,\ k=0,1\big\}=\left({\rm span}\{v,v_1\}\right)^{\bot}$,
\item[{\rm(ii)}]$S_1\ell^2(\Z):=\big\{f\in\ell^2(\Z):\left<f,v_{k}\right>=0,\ k=0,1,\ S_0T_0f=0\big\}$,
\item[{\rm(iii)}] $S_2\ell^2(\Z):=\big\{f\in\ell^2(\Z):\left<f,v_{k}\right>=0,\ k=0,1,2,\ QT_0f=0\big\}$,
\item[{\rm(iv)}] $S_3\ell^2(\Z):=\big\{f\in\ell^2(\Z):\left<f,v_{k}\right>=0,\ k=0,1,2,3,\ T_0f=0\big\}$.
\end{itemize}
\end{definition}
\begin{definition}\label{definition of Sjtilde}
Let $\widetilde{Q}=I-\widetilde{P}$, $\tilde{v}_k=Jv_k,k\in\N$ and 
$$\widetilde{G}_{2}(n,m)=-\frac{1}{24}|n-m|^3+\frac{5}{48}|n-m|-\big(2\sqrt2-3\big)^{|n-m|}\Big(\frac{1}{64}|n-m|+\frac{7\sqrt{2}}{256}\Big).$$
For any $j\in\{0,1,2\}$, define $\widetilde{S}_j$ as the orthogonal projection onto the space $\widetilde{S}_j\ell^2(\Z)$ with
\begin{itemize}
\item[{\rm(i)}]$\widetilde{S}_0\ell^2(\Z):=\big\{f\in\ell^2(\Z):\left<f,\tilde{v}\right>=0,\ \widetilde{Q}\widetilde{T}_0f=0\big\}$,
\item[{\rm(ii)}]$\widetilde{S}_1\ell^2(\Z):=\big\{f\in\ell^2(\Z):\left<f,\tilde{v}_k\right>=0,\ k=0,1,\ \widetilde{T}_0f=0\big\}$,
\item[{\rm(iii)}]$\widetilde{S}_2\ell^2(\Z):=\big\{f\in\ell^2(\Z):\left<f,\tilde{v}_k\right>=0,\ k=0,1,\ \widetilde{T}_0f=0,\ \widetilde{S}_1\tilde{v}\widetilde{G}_2\tilde{v}\widetilde{S}_1f=0\big\}$.
\end{itemize}
\end{definition}

\setcounter{section}{1}
\setcounter{equation}{19}
\setcounter{definition}{6}
\begin{theorem}\label{relation space and resonance types}
 Let $H=\Delta^2+V$ with $|V(n)|\lesssim \left<n\right>^{-\beta}$ for $\beta>9$. Then we have
\begin{itemize}
\item[{\rm(i)}] $0$ is a regular point of $H$ if and only if $S_1\ell^2(\Z)=\{0\}$.
\item[{\rm(ii)}] $0$ is a first kind resonance of $H$ if and only if $S_1\ell^2(\Z)\neq\{0\}$ and $S_2\ell^2(\Z)=\{0\}$.
\item[{\rm(iii)}] $0$ is a second kind resonance of $H$ if and only if $S_2\ell^2(\Z)\neq\{0\}$.

\item[{\rm(iv)}] $16$ is a regular point of $H$ if and only if $\widetilde{S}_0\ell^2(\Z)=\{0\}$.
\item[{\rm(v)}] $16$ is a resonance of $H$ if and only if $\widetilde{S}_0\ell^2(\Z)\neq\{0\}$ and $\widetilde{S}_{1}\ell^2(\Z)=\{0\}$.
\item[{\rm(vi)}] $16$ is an eigenvalue of $H$ if and only if $\widetilde{S}_{1}\ell^2(\Z)\neq\{0\}$.
\end{itemize}
\end{theorem}

\setcounter{section}{1}
\setcounter{equation}{19}
\setcounter{definition}{7}
\begin{theorem}\label{Asymptotic expansion theorem}
Let $H=\Delta^2+V$ with $|V(n)|\lesssim \left<n\right>^{-\beta}$ for some $\beta>0$. 
Then, given a sufficiently small $0<\mu_0\ll1$, the following asymptotic expansions hold in $\B(\ell^2(\Z))$ for any $0<\mu<\mu_0$.
\begin{itemize}
\item[{\rm(i)}] If $0$ is a regular point of $H$ and $\beta>15$, then
\begin{align}\label{asy expan of regular 0}
\left(M^{\pm}(\mu)\right)^{-1}=S_0A_{0}S_0+\mu QA^{\pm,0}_{11}Q+\mu^2(QA^{\pm,0}_{21}Q+S_0A^{\pm,0}_{22}+A^{\pm,0}_{23}S_0)+\mu^3A^{\pm,0}_3+\Gamma^{0}_{4}(\mu),
\end{align}
\item [{\rm(ii)}] If $0$ is a first kind resonance of $H$ and $\beta>19$, then
\begin{align}\label{asy expan of 1st 0}
\begin{split}
\left(M^{\pm}(\mu)\right)^{-1}&=\mu^{-1}S_{1}A^{\pm}_{-1}S_{1}+\big(S_0A^{\pm,1}_{01}Q+QA^{\pm,1}_{02}S_0\big)+\mu\big(S_0A^{\pm,1}_{11}+A^{\pm,1}_{12}S_0+QA^{\pm,1}_{13}Q\big)\\
&\quad+\mu^2\big(QA^{\pm,1}_{21}+A^{\pm,1}_{22}Q\big)+\mu^3A^{\pm,1}_{3}+\Gamma^{1}_{4}(\mu),
\end{split}
\end{align}
%\vskip0.2cm
\item [{\rm(iii)}] If $0$ is a second kind resonance of $H$ and $\beta>27$, then
\begin{align}\label{asy expan 2nd 0}
\left(M^{\pm}(\mu)\right)^{-1}&=\frac{S_{2}A^{\pm}_{-3}S_{2}}{\mu^3}+\frac{S_2A^{\pm}_{-2,1}S_0+S_0A^{\pm}_{-2,2}S_2}{\mu^2}
+\frac{S_2A^{\pm}_{-1,1}Q+QA^{\pm}_{-1,2}S_2+S_0A^{\pm}_{-1,3}S_0}{\mu}\notag\\
 &\quad+\big(S_2A^{\pm,2}_{01}+A^{\pm,2}_{02}S_2+QA^{\pm,2}_{03}S_0+S_0A^{\pm,2}_{04}Q\big)+\mu\big(S_0A^{\pm,2}_{11}+A^{\pm,2}_{12}S_0+QA^{\pm,2}_{13}Q\big)\notag\\
&\quad+\mu^2\big(QA^{\pm,2}_{21}+A^{\pm,2}_{22}Q\big)+\mu^3A^{\pm,2}_{3}+\Gamma^{2}_{4}(\mu),
\end{align}
\item[{\rm(iv)}] If $16$ is a regular point of $H$ and $\beta>7$, then
\begin{align}\label{asy expan regular 2}
\left(M^{\pm}(2-\mu)\right)^{-1}=\widetilde{Q}B_{0}\widetilde{Q}+
\mu^{\frac{1}{2}}B^{\pm,0}_1 +\Gamma^{0}_{1}(\mu),
\end{align}
\item [{\rm(v)}] If $16$ is a resonance of $H$ and $\beta>11$, then
\begin{align}\label{asy expan resonance 2}
\left(M^{\pm}(2-\mu)\right)^{-1}=\mu^{-\frac{1}{2}}\widetilde{S}_0B^{\pm}_{-1}\widetilde{S}_0+\big(\widetilde{S}_0B^{\pm,1}_{01}+B^{\pm,1}_{02}\widetilde{S}_0+\widetilde{Q}B^{\pm,1}_{03}\widetilde{Q}\big)+
\mu^{\frac{1}{2}}B^{\pm,1}_1 +\Gamma^{1}_{1}(\mu),
\end{align}
\item [{\rm(vi)}] If $16$ is an eigenvalue of $H$ and $\beta>15$, then
\begin{align}\label{asy expan eigenvalue 2}
\left(M^{\pm}(2-\mu)\right)^{-1}&=\mu^{-1}\widetilde{S}_1B_{-2}\widetilde{S}_1+\mu^{-\frac{1}{2}}\big(\widetilde{S}_0B^{\pm}_{-1,1}\widetilde{Q}+\widetilde{Q}B^{\pm}_{-1,2}\widetilde{S}_0\big)\notag\\
&\quad +\big(\widetilde{Q}B^{\pm,2}_{01}+B^{\pm,2}_{02}\widetilde{Q}\big)+
\mu^{\frac{1}{2}}B^{\pm,2}_1 +\Gamma^{2}_{1}(\mu),
\end{align}
\end{itemize}
where
$A_{0},A^{\pm}_{-1},A^{\pm}_{-3},A^{\pm}_{j,k},A^{\pm,i}_{jk},A^{\pm,i}_3,B_{0},B^{\pm}_{-1},B_{-2},B^{\pm}_{j,k},B^{\pm,i}_{jk},B^{\pm,i}_{1}$ are $\mu$-independent bounded operators on $\ell^2(\Z)$, Moreover, all the operators appearing on the right-hand
sides of \eqref{asy expan of regular 0}--\eqref{asy expan eigenvalue 2}
are absolutely bounded. The remainder terms $\Gamma^{i}_\ell(\mu)$ satisfy 
\begin{equation}\label{estimate of Gamma}
\|\Gamma^{i}_{\ell}(\mu)\|_{\ell^2\rightarrow\ell^2}+\mu\big\|\partial_{\mu}(\Gamma^{i}_{\ell}(\mu))\big\|_{\ell^2\rightarrow\ell^2}\lesssim\mu^{\ell},\qquad 0<\mu<\mu_0.
\end{equation}
\end{theorem}

\setcounter{section}{1}
\section{Limiting absorption principle}\label{sec of LAP}
\subsection{Free resolvent}\label{Subsec of free resol}
This subsection presents some basic properties of the  bi-Laplacian $\Delta^2$ on $\Z$. 
Recalling the definition of $\Delta$ in \eqref{definition of Laplacian}, the bi-Laplacian $\Delta^2$ is given by 
\begin{equation*}
(\Delta^2\phi)(n)=(\Delta(\Delta\phi))(n)=\phi(n+2)-4\phi(n+1)+6\phi(n)-4\phi(n-1)+\phi(n-2).
\end{equation*}
Consider the Fourier transform $\mcaF$: $\ell^2(\Z)\rightarrow L^2(\T)$ defined by
\begin{equation}\label{fourier transform}
(\mcaF\phi)(x):=(2\pi)^{-{\frac{1}{2}}}\sum_{n\in\Z}^{}e^{-inx}\phi(n), \quad x\in\T=\R/(2\pi\Z)\simeq(-\pi,\pi],\quad  \phi\in\ell^2(\Z).
\end{equation}
Under the Fourier transform, we have
\begin{equation}\label{unitary equivalent}
(\mcaF\Delta^2\phi)(x)=\mcaM(x)(\mcaF\phi)(x),\qquad \mcaM(x)=(2-2 \cos x)^2.
\end{equation}
It follows that the spectrum of $\Delta^2$ is purely absolutely
continuous and is given by $[0,16]$.
Let
  $$R_0(z):=(\Delta^2-z)^{-1},\qquad z\in\C\setminus[0,16],$$
  be the resolvent of $\Delta^2$, and define its boundary values on $(0,16)$ by
  \begin{align*}
  R^{\pm}_{0}(\lambda):=\lim\limits_{\varepsilon>0,\ \varepsilon\rightarrow0 }R_{0}(\lambda\pm i\varepsilon ),\qquad\lambda\in(0,16).
  \end{align*}
   Denote by $\B(s,s')$ the space of all bounded linear operators from $\ell^{2,s}(\Z)$ to $\ell^{2,s'}(\Z)$. The existence of $R^{\pm}_0(\lambda)$ in  $\B(s,-s)$ for $s>\frac{1}{2}$ follows from the resolvent decomposition
\begin{equation}\label{unity partition}
 R_0(z)=\frac{1}{2\sqrt z}\left(R_{-\Delta}(\sqrt z)-R_{-\Delta}(-\sqrt z)\right),\quad\sqrt{z}=\sqrt{|z|}e^{i\frac{\arg z}{2}},\quad  0<\arg z<2\pi,
 \end{equation}
together with the limiting absorption principle for $-\Delta$~(see, e.g., \cite[Lemma 3.1]{KKK06})
  $$R^{\pm}_{-\Delta}(\mu):=\lim\limits_{\varepsilon>0,\ \varepsilon\rightarrow0}R_{-\Delta}(\mu\pm i\varepsilon )\quad {\rm in\ the\ operator\ norm\ of }\ \B(s,-s)\ {\rm for}\ s>\frac{1}{2},\quad \mu\in(0,4),$$
where $R_{-\Delta}(\omega)=(-\Delta-\omega)^{-1}$ is the resolvent of $-\Delta$.

\setcounter{section}{5}
\setcounter{equation}{0}
\setcounter{definition}{0}
\begin{lemma}\label{Puiseux expan of R0}
Let $N$ be an integer and $\mu\in(0,2)$. 
\vskip0.1cm
{\rm (i)} For any $N\geq-3$, the following expansion holds in the norm of $\B(s,-s)$ for $s>\frac{9}{2}+N$:
\begin{equation}\label{Puiseux expan of R0 0}
R^{\pm}_0(\mu^4)=\sum\limits_{j=-3}^{N}\mu^{j}G^{\pm}_j+r^{\pm}_{N+1}(\mu),\qquad \left\|r^{\pm}_{N+1}(\mu)\right\|_{\B(s,-s)}=O\big(\mu^{N+1}\big),\quad\mu\rightarrow0^+,
\end{equation}
where $G^{\pm}_j\in\B(\sigma,-\sigma)$ for $\sigma>\frac{7}{2}+j$ with kernel $G^{\pm}_j(n-m)$ given by 
\begin{equation*}
   G^{\pm}_{2q-3}(n)= \frac{(\pm i(-1)^q-1)}{2^{2q+2}(2q)!}\prod\limits_{j=0}^{q-1}(4|n|^2-(2j+1)^2),\qquad  G^\pm_{2q-2}(n)=\frac{|n|(1+(-1)^{q+1})}{4(2q+1)!}\prod\limits_{j=1}^{q}(|n|^2-j^2).
\end{equation*}
In particular, $G^{\pm}_{4q+2}(n)=0$ for any $q\geq-1$, $G^{\pm}_{-3}(n)=\frac{-1\pm i}{4}$ and 
\begin{align}\label{expan coeffie of Gpm}
\begin{split}
  G^{\pm}_{-1}(n)&=\frac{1\pm i}{4}\big(\frac{1}{8}-\frac{1}{2}|n|^2\big), \qquad G^{\pm}_{1}(n)=\frac{-1\pm i}{32}\big(\frac{1}{3}|n|^4-\frac{5}{6}|n|^2+\frac{3}{16}\big), \\
 G^{\pm}_{0}(n)&=\frac{1}{12}\left(|n|^3-|n|\right), \qquad G^{\pm}_{3}(n)=\frac{-1\mp i}{4\cdot6!}\big(|n|^6-\frac{35}{4}|n|^4+\frac{259}{16}|n|^2-\frac{225}{64}\big).
\end{split}
\end{align}
Moreover, in the same sense, \eqref{Puiseux expan of R0 0} can be differentiated $N+4$ times in $\mu$.
\vskip0.2cm
{\rm(ii)}
For any $N\geq-1$, the following expansion holds in the norm of $\B(s,-s)$ for $s>\frac{5}{2}+N$:
\begin{equation}\label{Puiseux expan of R0 2}
JR^{\pm}_0((2-\mu)^4)J=\sum\limits_{j=-1}^{N}\mu^{\frac{j}{2}}\widetilde{G}^{\pm}_j+\mcaE^{\pm}_{N+1}(\mu),\qquad \left\|\mcaE^{\pm}_{N+1}(\mu)\right\|_{\B(s,-s)}=O\big(\mu^{\frac{N+1}{2}}\big),\quad \mu\rightarrow0^+,
\end{equation}
where $\widetilde{G}^{\pm}_j\in\B(\sigma,-\sigma)$ for $\sigma>\frac{3}{2}+j$ with kernel $\widetilde{G}^{\pm}_j(n-m)$. In particular, $\widetilde{G}^{\pm}_{-1}(n)=\frac{\pm i}{32}$ and
\begin{align}\label{expan coeffie of Gtutapm}
\begin{split}
\widetilde{G}^{\pm}_{0}(n)&=\frac{1}{32\sqrt2}\big(2\sqrt2|n|-{(2\sqrt2-3)}^{|n|}\big),\qquad  \widetilde{G}^{\pm}_{1}(n)=\frac{\mp i}{32}\big(2|n|^2-\frac{13}{8}\big),\\
\widetilde{G}^{\pm}_{2}(n)&=-\frac{1}{24}|n|^3+\frac{5}{48}|n|-\frac{\big(2\sqrt2-3\big)^{|n|}}{32}\Big(\frac{1}{2}|n|+\frac{7}{4\sqrt2}\Big).
\end{split}
\end{align}
Furthermore, in the same sense, \eqref{Puiseux expan of R0 2} can be differentiated $N+2$ times in $\mu$.
\end{lemma}
\begin{proof}
Let $x=|n|$ and fix $0<\mu_0\ll1$. Denote $\Phi^{\pm}(\mu,x)=\pm i a_1(\mu)e^{\mp ix\theta(\mu)}+a_2(\mu)e^{b(\mu)x}$. 
\vskip0.1cm
\textbf{(i) Expansion near $\mu=0$.}
For any integer $N\geq-3$, expanding $\Phi^{\pm}(\mu,x)$ in a Taylor series at $\mu=0$ up to order $N+3$, we obtain
\[
R^{\pm}_0(\mu^4,n)=\frac{1}{4}\mu^{-3}\Phi^{\pm}(\mu,x)
=\sum_{j=-3}^{N}\mu^jG_j^{\pm}(x)+r_{N+1}^{\pm}(\mu,x),
\]
where
\[
G_j^{\pm}(x)=\frac{(\partial_\mu^{j+3}\Phi^{\pm})(0,x)}{4(j+3)!},\qquad
r_{N+1}^{\pm}(\mu,x)
=\frac{\mu^{N+1}}{4(N+3)!}\int_0^1(1-t)^{N+3}
(\partial_\mu^{N+4}\Phi^{\pm})(t\mu,x)\,dt.
\]
Since $a_1,a_2,\theta,b\in C^{\infty}([0,\mu_0])$, $b(0)=0$, and $b(\mu)<0$ for $\mu>0$, we have, for $0\leq k\leq N+4$,
\[
\bigl|(\partial_\mu^k\Phi^{\pm})(\mu,x)\bigr|\lesssim\left<x\right>^k,
\qquad 0\leq\mu\leq\mu_0.
\]
Indeed, $|e^{\mp ix\theta(\mu)}|=1$ and $|e^{b(\mu)x}|\leq1$, and each $\mu$-derivative produces at most one additional factor of $x$. Differentiating the Taylor remainder (equivalently, applying Taylor's formula to the corresponding derivatives) and then using the Leibniz rule for the factor $\mu^{-3}$ yields
\begin{equation}\label{estim for remainder r}
\bigl|(\partial_\mu^{\ell}r_{N+1}^{\pm})(\mu,x)\bigr|
\lesssim \mu^{N+1-\ell}\left<x\right>^{N+4},
\qquad 0<\mu<\mu_0,\quad 0\leq\ell\leq N+4.  
\end{equation}

To compute $G_j^{\pm}(x)$, recall that $\theta(\mu)=-2\arcsin\frac{\mu}{2}$ and $b(\mu)=-2\operatorname{arsinh}\frac{\mu}{2}$. We introduce 
$$\Phi^{\pm}_1(\mu,x)=e^{\pm 2ix\arcsin{\mu}},\qquad \Phi_2(\mu,x)=e^{-2x\operatorname{arsinh}\mu}.$$ 
For each fixed $x$, these functions satisfy
\begin{align}\label{rel Phipm1 Phi2}
\begin{split}
(1-\mu^2)(\partial^2_\mu\Phi^{\pm}_1)(\mu,x)+4x^2\Phi^{\pm}_1(\mu,x)&=\mu (\partial_\mu\Phi^{\pm}_1)(\mu,x),\\
(1+\mu^2)(\partial^2_\mu\Phi_2)(\mu,x)+\mu(\partial_\mu\Phi_2)(\mu,x)&=4x^2\Phi_2(\mu,x).
\end{split}
\end{align}
Write
\begin{equation}\label{talor expa of Phi1 2}
\Phi^{\pm}_1(\mu,x)=\sum\limits_{k=0}^{\infty}C^{\pm}_k(x)\mu^k,\qquad \Phi_2(\mu,x)=\sum\limits_{k=0}^{\infty}D_k(x)\mu^k,\qquad 0<\mu<\mu_0.
\end{equation}
In particular, $C^{\pm}_0(x)=D_0(x)=1$, $C^{\pm}_1(x)=\pm 2ix$ and $D_1(x)=-2x$. Substituting the series into  \eqref{rel Phipm1 Phi2} gives
\[
C_{k+2}^{\pm}(x)=\frac{k^2-4x^2}{(k+2)(k+1)}C_k^{\pm}(x),
\qquad
D_{k+2}(x)=\frac{4x^2-k^2}{(k+2)(k+1)}D_k(x).
\]
Consequently,
\[
C_{2q}^{\pm}(x)=(-1)^qD_{2q}(x),
\qquad
C_{2q+1}^{\pm}(x)=\mp i(-1)^qD_{2q+1}(x),
\qquad q\geq0,
\]
where, with the convention that an empty product equals $1$,
\begin{equation}\label{Dkx}
   D_{2q}(x)=
       \frac{1}{(2q)!}\prod\limits_{j=0}^{q-1}(4x^2-4j^2),\qquad D_{2q+1}(x)=
       \frac{-2x}{(2q+1)!}\prod\limits_{j=0}^{q-1}(4x^2-(2j+1)^2).
\end{equation}
Since $\Phi^{\pm}(\mu,0)=\pm ia_1(\mu)+a_2(\mu)$, and for $x\neq0$, 
$$\Phi^{\pm}(\mu,x)=
    \frac{1}{2x}\big((\partial_{\mu}\Phi^{\pm}_1)(\frac{\mu}{2},x)+(\partial_{\mu}\Phi_2)(\frac\mu2,x)\big).$$ 
Combining the preceding identities with \eqref{talor expa of Phi1 2} and \eqref{Dkx}, we obtain, for all $x\geq0$,
\begin{equation*}
R^{\pm}_0(\mu^4,x)=\sum\limits_{q=0}^{+\infty} G^{\pm}_{2q-3}(x)\mu^{2q-3}+\sum\limits_{q=0}^{+\infty} G^\pm_{2q-2}(x)\mu^{2q-2}, \qquad 0<\mu\leq \mu_0,
\end{equation*}
where
\begin{equation*}
   G^{\pm}_{2q-3}(x)= \frac{(\pm i(-1)^q-1)}{2^{2q+2}(2q)!}\prod\limits_{j=0}^{q-1}(4x^2-(2j+1)^2),\qquad  G^\pm_{2q-2}(x)=\frac{x(1+(-1)^{q+1})}{4(2q+1)!}\prod\limits_{j=1}^{q}(x^2-j^2).
\end{equation*}
In particular, $G^\pm_{4k+2}(x)=0$ for $k\geq-1$. Since $\left<n-m\right>^{2j}\lesssim\left<n\right>^{2j}\left<m\right>^{2j}$, it follows that $$\sum\limits_{n,m\in\Z}\left<n\right>^{-2s}\left<n-m\right>^{2j}\left<m\right>^{-2s}<\infty$$ for any $s>1/2+j$. This, together with \eqref{estim for remainder r} proves \eqref{Puiseux expan of R0 0}.

\vskip0.1cm
\textbf{(ii) Expansion near the upper endpoint.}
Since $(-1)^{n+m}=(-1)^{|n-m|}$, we have
\[
\bigl(JR_0^{\pm}((2-\mu)^4)J\bigr)(n,m)
=(-1)^{|n-m|}R_0^{\pm}((2-\mu)^4,n-m).
\]
Hence it suffices to study, with $x\in\N\cup\{0\}$,
\begin{equation*}
\widetilde{R}^{\pm}_0(\mu,x)=\frac{(-1)^x}{4(2-\mu)^3}\Phi^{\pm}(2-\mu,x).
\end{equation*}
By the definition of $\theta(2-\mu)$, it follows that
$\theta(2-\mu)=-\pi+\alpha(\mu)$ with $0<\alpha(\mu)=4\arcsin{\frac{\sqrt{\mu}}{2}}.$
Let $T_x$ and $U_x$ denote the Chebyshev polynomials of the first and second kind, respectively. Thus
\[
T_x(y)=
\begin{cases}
\frac{x}{2}\sum\limits_{k=0}^{[x/2]}(-1)^k
\frac{(x-k-1)!}{k!(x-2k)!}(2y)^{x-2k},&x\geq1,\\
1,&x=0,
\end{cases},\ \ 
U_x(y)=\sum_{k=0}^{[x/2]}(-1)^k
\binom{x-k}{k}(2y)^{x-2k}\ {\rm for}\  x\geq0,
\]
with the convention $U_{-1}=0$.
Using Euler's formula, $\cos\alpha(\mu)=1-2\mu+\frac{\mu^2}{2}$ and 
$$\cos (\alpha x)=T_x(\cos \alpha),\qquad \sin (\alpha x)= U_{x-1}(\cos\alpha)\sin \alpha,$$
we obtain
\begin{equation}\label{eq:upper-cheb}
\widetilde{R}^{\pm}_0(\mu,x)=\frac{\pm iT_x(1+\beta(\mu))}{2\sqrt{\mu}(2-\mu)^3\sqrt{(4-\mu)}}+\frac{U_{x-1}(1+\beta(\mu))}{4(2-\mu)^2}-\frac{(q(\mu))^x}{2(2-\mu)^3d(\mu)}
\end{equation}
where
$$\beta(\mu)=\frac{\mu^2}{2}-2\mu,\qquad d(\mu)=\sqrt{8-4\mu+\mu^2},\qquad q(\mu)=-\frac{1}{4}\big(2-\mu-d(\mu)\big)^2.$$
Since $T_x$ has degree $x$ for $x\geq0$ and $U_{x-1}$ has degree $x-1$ for $x\geq1$ (with $U_{-1}=0$), we have
\[
T_x(1+\beta(\mu))
=\sum_{\ell=0}^{x}\frac{T_x^{(\ell)}(1)}{\ell!}(\beta(\mu))^\ell
,\qquad
U_{x-1}(1+\beta(\mu))
=\sum_{\ell=0}^{x-1}\frac{U_{x-1}^{(\ell)}(1)}{\ell!}(\beta(\mu))^\ell,
\]
where $T_x(1)=1$, $U_{x-1}(1)=x$, and for $1\leq j\leq x $ and $1\leq k\leq x-1$,
\begin{equation}\label{Tx1 Ux1}
T_x^{(j)}(1)=
    \frac{x^2 P_{j-1}(x)}
{(2j-1)!!}
,\qquad U_{x-1}^{(k)}(1)=
    \frac{xP_{k}(x)}
{(2k+1)!!}~
,\qquad P_{\ell}(x)=\displaystyle\prod_{s=1}^{\ell}(x^2-s^2).
\end{equation}
Write $q(\mu)=q(0)+\mu h(\mu)$, where $q(0)=2\sqrt{2}-3$ and  $h(\mu)$ is analytic in $\mu$. It follows that $(q(\mu))^x$ is analytic in $\mu$ and contains only integer powers
of $\mu$. Moreover, there exists $\rho\in(0,1)$ such that
\begin{equation}\label{qu}
\left|\partial_\mu^r (q(\mu))^x\right|
\lesssim_r \left<x\right>^r\rho^x,
\qquad 0\leq\mu\leq\mu_0,\quad r\in\N.
\end{equation}
Thus \eqref{eq:upper-cheb} becomes
\begin{equation}\label{final expan}
\widetilde R_0^\pm(\mu,x)
=
\sum_{k=0}^{\infty}
\mu^{k-\frac12}\widetilde G_{2k-1}^\pm(x)
+
\sum_{k=0}^{\infty}
\mu^k\widetilde G_{2k}(x),\qquad0<\mu<\mu_0, 
\end{equation} 
where $\widetilde G_{2k-1}^\pm(x)=\frac{\pm i}{2}E_k(x)$, $\widetilde G_{2k}(x)=F_k(x)$, and $E_k(x)$ and $F_k(x)$ are given by
\begin{equation*}
E(\mu,x)=\sum_{\ell=0}^{x}
\frac{T_x^{(\ell)}(1)(\beta(\mu))^\ell}
{(2-\mu)^3\sqrt{4-\mu}\,\ell!}=\sum_{k=0}^\infty E_k(x)\mu^k,\qquad0<\mu<\mu_0
\end{equation*}
and 
\begin{equation*}
F(\mu,x)=\sum_{\ell=0}^{x-1}
\frac{U_{x-1}^{(\ell)}(1)(\beta(\mu))^\ell}
{4(2-\mu)^2\ell!}
-\frac{(q(\mu))^x}{2(2-\mu)^3d(\mu)}=\sum\limits_{k=0}^{\infty}F_k(x)\mu^k,\qquad0<\mu<\mu_0.
\end{equation*}
The formulas \eqref{Tx1 Ux1} and \eqref{qu} above imply
$
E_k(x)=O(\left< x\right>^{2k})$,
while $F_k(x)$ is the sum of a polynomial-growth term of order at most
$2k+1$ and an exponentially decaying term. More generally, for any integer $N\geq-1$, truncating
\eqref{final expan} gives the desired \eqref{Puiseux expan of R0 2}. This completes the proof.
\end{proof}

\setcounter{section}{5}
\setcounter{equation}{12}
\setcounter{definition}{1}
\begin{lemma}\label{lemma of kernel and operaor}
Let $G_0=G^{\pm}_0$,  $\widetilde{G}_0=\widetilde{G}^{\pm}_0$, $\widetilde{G}_2=\widetilde{G}^{\pm}_2$, and $$G_{-1}=\frac{4}{1\pm i}G^{\pm}_{-1},\quad G_{1}=\frac{32}{-1\pm i}G^{\pm}_{1},\quad G_3=\frac{4\cdot6!}{-1\mp i}G^{\pm}_{3}(n),\quad \widetilde{G}_1=\mp i\widetilde{G}^{\pm}_1.$$
 Let $T_0=U+vG_0v$ and $\widetilde{T}_0=U+\tilde{v}\widetilde{G}_0\tilde{v}$. Then the following statements hold.
\vskip0.1cm
{\rm(1)}$~
{\rm Ker}QvG_{-1}vQ\big|_{Q\ell^2(\Z)}=S_0\ell^2(\Z).$
\vskip0.2cm
{\rm(2)} Denote by $D_0$ the inverse of $QvG_{-1}vQ+S_0$ on $Q\ell^2(\Z)$ and define
 \begin{equation}\label{T1}
T_1=S_1vG_1vS_1+\frac{8}{\|V\|_{\ell^1}}S_1vG_{-1}vPvG_{-1}vS_1+64S_1T_0D_0T_0S_1.
\end{equation}
Then
${\rm Ker}T_1\big|_{S_1\ell^2(\Z)}=S_2\ell^2(\Z).$
\vskip0.2cm
{\rm(3)} Denote by $D_2$ the inverse of $T_1+S_2$ on $S_1\ell^2(\Z)$ and let
\begin{align}\label{T2}
T_2&=\frac{1}{6!}\Big(S_2vG_3vS_2-\frac{8\cdot6!}{\|V\|_{\ell^1}}S_2T^2_0S_2-\frac{6!}{64}S_2vG_1vD_0vG_1vS_2\Big)\notag\\
&\quad+\frac{64}{\|V\|^2_{\ell^1}}\big(S_2T_0vG_{-1}vD_0-\frac{\|V\|_{\ell^1}}{8}S_2vG_1vD_0T_0D_0\big)D_2\\
&\quad\times \big(D_0vG_{-1}vT_0S_2-\frac{\|V\|_{\ell^1}}{8}D_0T_0D_0vG_1vS_2\big).\notag
\end{align}
Then
${\rm Ker}T_2\big|_{S_2\ell^2(\Z)}=S_3\ell^2(\Z).$
\vskip0.2cm
{\rm(4)} Define
\begin{equation}\label{T1tuta}
\widetilde{T}_1=\widetilde{S}_0\tilde{v}\widetilde{G}_1\tilde{v}\widetilde{S}_0+\frac{32}{\|V\|_{\ell^1}}\widetilde{S}_0\widetilde{T}^2_0\widetilde{S}_0.
\end{equation}
Then
${\rm Ker}\widetilde{T}_1\big|_{\widetilde{S}_0\ell^2(\Z)}=\widetilde{S}_1\ell^2(\Z).$
\vskip0.2cm
{\rm(5)} Let $H=\Delta^2+V$ on $\Z$ and $|V(n)|\lesssim\left<n\right>^{-\beta}$ with $\beta>9$. Define $\widetilde{T}_2=\widetilde{S}_1\tilde{v}\widetilde{G}_2\tilde{v}\widetilde{S}_1$. Then
$$S_3\ell^2(\Z)=\{0\},\qquad{\rm Ker}\widetilde{T}_2\big|_{\widetilde{S}_1\ell^2(\Z)}=\{0\}.$$
\end{lemma}
\begin{remark}\label{remark of charac from inverti of operators}
{\rm Using Theorem \ref{relation space and resonance types} and 
$$S_1\ell^2(\Z)={\rm Ker}S_0T_0S_0\big|_{S_0\ell^2(\Z)},\qquad \widetilde{S}_0\ell^2(\Z)={\rm Ker}\widetilde{Q}\widetilde{T}_0\widetilde{Q}\big|_{\widetilde{Q}\ell^2(\Z)},$$
the lemma further yields the following characterizations of the resonance types of $H$:
\vskip0.1cm
(i) 0 is a regular point of $H$ if and only if $S_0T_0S_0$ is invertible on $S_0\ell^2(\Z)$.

(ii) 0 is a first kind resonance of $H$ if and only if $S_0T_0S_0$ is not invertible on $S_0\ell^2(\Z)$ but $T_1$ is invertible on $S_1\ell^2(\Z)$.

(iii) 0 is a second kind resonance of $H$ if and only if $T_1$ is not invertible on $S_1\ell^2(\Z)$.

(iv) 16 is a regular point of $H$ if and only if $\widetilde{Q}\widetilde{T}_0\widetilde{Q}$ is invertible on $\widetilde{Q}\ell^2(\Z)$.

(v) 16 is a resonance of $H$ if and only if $\widetilde{Q}\widetilde{T}_0\widetilde{Q}$ is not invertible on $\widetilde{Q}\ell^2(\Z)$ but $\widetilde{T}_1$ is invertible on $\widetilde{S}_0\ell^2(\Z)$.

(vi) 16 is an eigenvalue of $H$ if and only if $\widetilde{T}_1$ is not invertible on $\widetilde{S}_0\ell^2(\Z)$.
}
\end{remark}
The proof of this lemma will be delayed to the end of this section. Furthermore, the following lemma will be frequently used in the proof.
\begin{lemma}\label{lemm of expa}
{\rm (\cite[Lemma 2.1]{JN01})} Let $\mscH$ be a complex Hilbert space. Let $A$ be a closed operator and $S$ a projection. Suppose $A+S$ has a bounded inverse. Then $A$ has a bounded inverse if and only if
$$a\equiv S-S(A+S)^{-1}S$$
has a bounded inverse in $S\mscH$, and in this case
$$A^{-1}=(A+S)^{-1}+(A+S)^{-1}Sa^{-1}S(A+S)^{-1}.$$
\end{lemma}
With these three lemmas in hand, we now begin the proof of Theorem \ref{Asymptotic expansion theorem}.

\setcounter{section}{5}
\setcounter{equation}{15}
\setcounter{definition}{4}
\begin{proof}[Proof of Theorem {\rm\ref{Asymptotic expansion theorem}}]{\bf{(i) ${\bm{0}}$ is a regular point of $H$ and {$\bm{\beta>15}$}.}}
Taking $N=3$ in \eqref{Puiseux expan of R0 0}, we have the expansion in $\B(s,-s)$ for $s>\frac{15}{2}$,
\begin{equation*}
R^{\pm}_0(\mu^4)=\mu^{-3}G^{\pm}_{-3}+\mu^{-1}G^{\pm}_{-1}+G^{\pm}_0+\sum\limits_{k=1}^{3}\mu^kG^{\pm}_{k}+\Gamma_4(\mu),\qquad \mu \rightarrow0.
\end{equation*}
Since $\beta>15$ and $M^{\pm}(\mu)=U+vR^{\pm}_0(\mu^4)v$, using $vG^{\pm}_{-3}v=a^{\pm}P$ with $a^{\pm}=\frac{-1\pm i}{4}\|V\|_{\ell^1}$, we obtain 
 \begin{equation*}
 M^{\pm}(\mu)=\frac{a^{\pm}}{\mu^3}\Big(P+\frac{\mu^2}{a^{\pm}}vG^{\pm}_{-1}v+\frac{\mu^3}{a^{\pm}}T_0+\frac{1}{a^{\pm}}\sum\limits_{k=1}^{3}\mu^{k+3}vG^{\pm}_{k}v+\Gamma_7(\mu)\Big)=:\frac{a^{\pm}}{\mu^3}\widetilde{M}^{\pm}(\mu)
 \end{equation*}
 on $\ell^2(\Z)$ as $\mu\rightarrow0$. Then
\begin{equation}\label{Relat of M and M tuta} \left(M^{\pm}(\mu)\right)^{-1}=\frac{\mu^3}{a^{\pm}}\big(\widetilde{M}^{\pm}(\mu)\big)^{-1}.
\end{equation}

{\bf{\underline{Step1.}}} By Lemma \ref{lemm of expa}, $\widetilde{M}^{\pm}(\mu)$ is invertible on $\ell^{2}(\Z)$ if and only if $$ M^{\pm}_1(\mu):=Q-Q\big(\widetilde{M}^{\pm}(\mu)+Q\big)^{-1}Q$$ is invertible on $Q\ell^{2}(\Z)$ and in this case, one has
\begin{align}\label{Relat of M tuta and M1}
\big(\widetilde{M}^{\pm}(\mu)\big)^{-1}=\big(\widetilde{M}^{\pm}(\mu)+Q\big)^{-1}\Big[I+Q\left(M^{\pm}_1(\mu)\right)^{-1}Q\big(\widetilde{M}^{\pm}(\mu)+Q\big)^{-1}\Big].
\end{align}
Using Neumann series expansions, we obtain 
\begin{align}\label{expr of Bkpm}
\big(\widetilde{M}^{\pm}(\mu)+Q\big)^{-1}=I-\sum\limits_{k=1}^{5}\mu^{k+1}B^{\pm}_{k}+\Gamma_7(\mu),\quad \mu\rightarrow0,
\end{align}
where
\begin{itemize}
\item $B^{\pm}_{1}=\frac{1}{a^{\pm}}vG^{\pm}_{-1}v$,\quad$B^{\pm}_{2}=\frac{1}{a^{\pm}}T_0$,\quad $B^{\pm}_{3}=\frac{1}{a^{\pm}}vG^{\pm}_{1}v-\left(\frac{1}{a^{\pm}}vG^{\pm}_{-1}v\right)^2$,
\vskip0.2cm
\item $B^{\pm}_{4}=-\left(\frac{1}{a^{\pm}}\right)^2\left(vG^{\pm}_{-1}vT_0+T_0vG^{\pm}_{-1}v\right)$,
\item
$B^{\pm}_{5}=\frac{1}{a^{\pm}}vG^{\pm}_{3}v-\left(\frac{1}{a^{\pm}}\right)^2\left(vG^{\pm}_{-1}v\cdot vG^{\pm}_{1}v+T^2_0+vG^{\pm}_{1}v \cdot vG^{\pm}_{-1}v\right)+\left(\frac{1}{a^{\pm}}vG^{\pm}_{-1}v\right)^3$.
\end{itemize}
This, together with $Q^2=Q$, and $QB^{\pm}_1Q=b^{\pm}QvG_{-1}vQ$ with $b^{\pm}=\mp i\|V\|^{-1}_{\ell^1}$, gives
\begin{equation*}
M^{\pm}_1(\mu)=b^{\pm}\mu^2\Big(QvG_{-1}vQ+\frac{1}{b^{\pm}}\sum\limits_{k=2}^{5}\mu^{k-1}QB^{\pm}_{k}Q+Q\Gamma_5(\mu)Q\Big)=:b^{\pm}\mu^2\widetilde{M}^{\pm}_1(\mu),
\end{equation*}
Hence,
\begin{equation}\label{Relat of M1 and M1 tuta}
\left(M^{\pm}_1(\mu)\right)^{-1}=\frac{1}{b^{\pm}\mu^2}\big(\widetilde{M}_{1}^{\pm}(\mu)\big)^{-1}.
\end{equation}
\vskip0.15cm
{\bf{\underline{Step2.}}} Since ${\rm Ker}QvG_{-1}vQ\big|_{Q\ell^2(\Z)}=S_0\ell^{2}(\Z)$, Lemma \ref{lemm of expa} implies that $\widetilde{M}_{1}^{\pm}(\mu)$ is invertible on $Q\ell^{2}(\Z)$ if and only if $$M^{\pm}_{2}(\mu):=S_0-S_0\big(\widetilde{M}_{1}^{\pm}(\mu)+S_0\big)^{-1}S_0$$ is invertible on $S_0\ell^{2}(\Z)$. In this case,
\begin{equation}\label{Relat of M1 tuta and M2}
\big(\widetilde{M}_{1}^{\pm}(\mu)\big)^{-1}=\big(\widetilde{M}_{1}^{\pm}(\mu)+S_0\big)^{-1}\Big(I+S_0\big(M^{\pm}_{2}(\mu)\big)^{-1}S_0\big(\widetilde{M}_{1}^{\pm}(\mu)+S_0\big)^{-1}\Big),
\end{equation}
where
\begin{equation}\label{expre of Bpm tuta}
\big(\widetilde{M}_{1}^{\pm}(\mu)+S_0\big)^{-1}=D_0-\sum\limits_{k=1}^{4}\mu^{k}\widetilde{B}_{k}^{\pm}+D_0\Gamma_5(\mu)D_0,\qquad\mu\rightarrow0,
\end{equation}
with $D_0$ being the inverse of $QvG_{-1}vQ+S_0$ on $Q\ell^2(\Z)$~(hence $D_0Q=D_0=QD_0$) and
\begin{itemize}
\item $\widetilde{B}_{1}^{\pm}=\frac{1}{b^{\pm}}D_0B^{\pm}_2D_0$,\quad$\widetilde{B}_{2}^{\pm}=\frac{1}{b^{\pm}}D_0B^{\pm}_3D_0-\left(\frac{1}{b^{\pm}}\right)^2\left(D_0B^{\pm}_2\right)^2D_0$,
\vskip0.2cm
\item
$\widetilde{B}_{3}^{\pm}=\frac{1}{b^{\pm}}D_0B^{\pm}_4D_0-\left(\frac{1}{b^{\pm}}\right)^2\left(D_0B^{\pm}_2D_0B^{\pm}_3D_0+D_0B^{\pm}_3D_0B^{\pm}_2D_0\right)+\left(\frac{1}{b^{\pm}}D_0B^{\pm}_2\right)^3D_0,$
\vskip0.2cm
\item $\widetilde{B}_{4}^{\pm}=\frac{1}{b^{\pm}}D_0B^{\pm}_5D_0-\left(\frac{1}{b^{\pm}}\right)^2\left(D_0B^{\pm}_2D_0B^{\pm}_4D_0+\left(D_0B^{\pm}_3\right)^2D_0+D_0B^{\pm}_4D_0B^{\pm}_2D_0\right)$

$\quad+\left(\frac{1}{b^{\pm}}\right)^3\left(\left(D_0B^{\pm}_2\right)^2D_0B^{\pm}_3D_0+D_0B^{\pm}_2D_0B^{\pm}_3D_0B^{\pm}_2D_0+D_0B^{\pm}_3\left(D_0B^{\pm}_2\right)^2D_0\right)$

$\quad-\left(\frac{1}{b^{\pm}}D_0B^{\pm}_2\right)^4D_0$.
\end{itemize}
Using $S_0D_0=S_0=D_0S_0$ and $S_0\widetilde{B}_{1}^{\pm}S_0=\frac{1}{a^{\pm}_{-1}}S_0T_0S_0$ with $a^{\pm}_{-1}=\frac{1\pm i}{4}$, we obtain
\begin{align*}
M^{\pm}_2(\mu)=\frac{\mu}{a^{\pm}_{-1}}\Big(S_0T_0S_0+a^{\pm}_{-1}\sum\limits_{k=2}^{4}\mu^{k-1}S_0\widetilde{B}_{k}^{\pm}S_0+S_0\Gamma_4(\mu)S_0\Big)=:\frac{\mu}{a^{\pm}_{-1}}\widetilde{M}^{\pm}_2(\mu),
\end{align*}
and thus
\begin{equation}\label{Relat of M2 and M2 tuta}
\left(M^{\pm}_2(\mu)\right)^{-1}=\frac{a^{\pm}_{-1}}{\mu}\big(\widetilde{M}_{2}^{\pm}(\mu)\big)^{-1}.
\end{equation}
Since $0$ is a regular point of $H$, Remark \ref{remark of charac from inverti of operators} implies that $S_0T_0S_0$ is invertible on $S_0\ell^2(\Z)$. Denote its inverse by $D_1$. Then, using the Neumann expansion and $S_0D_1=D_1=D_1S_0$, we obtain
\begin{align}\label{inverse of Mpm2}
\left(M^{\pm}_2(\mu)\right)^{-1}=\frac{a^{\pm}_{-1}}{\mu}D_1+\sum\limits_{k=0}^{2}\widetilde{D}^{\pm}_k\mu^k+D_1\Gamma_3(\mu)D_1,\quad\mu\rightarrow0,
\end{align}
where
\begin{itemize}
\item
$\widetilde{D}^{\pm}_0=-\left(a^{\pm}_{-1}\right)^2D_1\widetilde{B}_{2}^{\pm}D_1$,\quad $\widetilde{D}^{\pm}_{1}=-\big(a^{\pm}_{-1}\big)^2\big(D_1\widetilde{B}_{3}^{\pm}D_1-a^{\pm}_{-1}\big(D_1\widetilde{B}_{2}^{\pm}\big)^2D_1\big)$,
\vskip0.2cm
\item $\widetilde{D}^{\pm}_2=-\big(a^{\pm}_{-1}\big)^2D_1\widetilde{B}_{4}^{\pm}D_1+\big(a^{\pm}_{-1}\big)^3\big(D_1\widetilde{B}_{2}^{\pm}D_1\widetilde{B}_{3}^{\pm}D_1+D_1\widetilde{B}_{3}^{\pm}D_1\widetilde{B}_{2}^{\pm}D_1\big)$
    $-\big(a^{\pm}_{-1}\big)^4\big(D_1\widetilde{B}_{2}^{\pm}\big)^3D_1$.
\end{itemize}
Taking \eqref{inverse of Mpm2} into \eqref{Relat of M1 tuta and M2},\eqref{Relat of M1 and M1 tuta},\eqref{Relat of M tuta and M1} and \eqref{Relat of M and M tuta} in order, the desired \eqref{asy expan of regular 0} is obtained.
\vskip0.2cm
{\bf{(ii) ${\bm{0}}$ is a first kind resonance of $H$ and {$\bm{\beta>19}$}.}} For this case, we take $N=5$ in \eqref{Puiseux expan of R0 0} and repeat {\bf{\underline{Step1}}} and  {\bf{\underline{Step2}}} in {\bf{(i)}}. However, since $S_0T_0S_0$ is not invertible on $S_0\ell^2(\Z)$, we need to further analyze the $\widetilde{M}^{\pm}_2(\mu)$ in {\bf{\underline{Step2}}} of {\bf{(i)}}. Since ${\rm Ker}S_0T_0S_0\big|_{S_0\ell^2(\Z)}=S_1\ell^2(\Z)$, Lemma \ref{lemm of expa} yields the following further reduction.
\vskip0.1cm
{\bf{\underline{Step3.}}} \ $\widetilde{M}_{2}^{\pm}(\mu)$ is invertible on $S_0\ell^{2}(\Z)$ if and only if $$ M^{\pm}_{3}(\mu):=S_1-S_1\big(\widetilde{M}_{2}^{\pm}(\mu)+S_1\big)^{-1}S_1$$ is invertible on $S_1\ell^{2}(\Z)$. In this case,
\begin{equation}\label{Relat of M2 tuta and M3}
\big(\widetilde{M}_{2}^{\pm}(\mu)\big)^{-1}=\big(\widetilde{M}_{2}^{\pm}(\mu)+S_1\big)^{-1}\Big(I+S_1\big(M^{\pm}_{3}(\mu)\big)^{-1}S_1\big(\widetilde{M}_{2}^{\pm}(\mu)+S_1\big)^{-1}\Big).
\end{equation}
Let $D_1$ be the inverse of $S_0T_0S_0+S_1$ on $S_0\ell^2(\Z)$.~Then $D_1S_0=D_1=S_0D_1$ and
\begin{equation}\label{expre of Bpm ttuta}
\big(\widetilde{M}_{2}^{\pm}(\mu)+S_1\big)^{-1}=D_1-\sum\limits_{k=1}^{5}\mu^{k}\tilde{\tilde{B}}^{\pm}_{k}+D_1\Gamma_6(\mu)D_1,\quad\mu\rightarrow0,
\end{equation}
where
\begin{itemize}
\item $\tilde{\tilde{B}}^{\pm}_{1}=a^{\pm}_{-1}D_1\widetilde{B}^{\pm}_2D_1$,\quad$\tilde{\tilde{B}}^{\pm}_{2}=a^{\pm}_{-1}D_1\widetilde{B}^{\pm}_3D_1-\big(a^{\pm}_{-1}D_1\widetilde{B}^{\pm}_2\big)^2D_1$,
\vskip0.2cm
\item $\tilde{\tilde{B}}^{\pm}_{3}=a^{\pm}_{-1}D_1\widetilde{B}^{\pm}_4D_1-{(a^{\pm}_{-1})}^2\big(D_1\widetilde{B}^{\pm}_2D_1\widetilde{B}^{\pm}_3D_1+D_1\widetilde{B}^{\pm}_3D_1\widetilde{B}^{\pm}_2D_1\big)+\big(a^{\pm}_{-1}D_1\widetilde{B}^{\pm}_2\big)^3D_1$.
\end{itemize}
Using $S_1D_1=S_1=D_1S_1$, $S_1D_0=S_1=D_0S_1$ and $S_1\tilde{\tilde{B}}^{\pm}_{1}S_1=a^{\pm}_{1}T_1$ with $a^{\pm}_{1}=\frac{-1\pm i}{32}$, we obtain
\begin{align*}
M^{\pm}_3(\mu)=a^{\pm}_{1}\mu\Big(T_1+\frac{1}{a^{\pm}_{1}}\sum\limits_{k=2}^{5}\mu^{k-1}S_1\tilde{\tilde{B}}^{\pm}_{k}S_1+S_1\Gamma_5(\mu)S_1\Big)=:a^{\pm}_{1}\mu\widetilde{M}^{\pm}_3(\mu),
\end{align*}
and then
\begin{equation}\label{Relat of M3 and M3 tuta}
\left(M^{\pm}_3(\mu)\right)^{-1}=\frac{1}{a^{\pm}_{1}\mu}\left(\widetilde{M}_{3}^{\pm}(\mu)\right)^{-1}.
\end{equation}
By Remark \ref{remark of charac from inverti of operators}, since $T_1$ is invertible on $S_1\ell^2(\Z)$, using Von-Neumann expansion, we derive
\begin{align}\label{inverse of Mtuta3}
\big(\widetilde{M}_{3}^{\pm}(\mu)\big)^{-1}
&=D_2-\sum\limits_{k=1}^{4}\mu^k\tilde{\tilde{\tilde{B}}}^{\pm}_{k}+D_2\Gamma_5(\mu)D_2,\quad \mu\rightarrow0,
\end{align}
where $D_2$ is the inverse of $T_1$ on $S_1\ell^2(\Z)$~(then $D_2S_1=D_2=S_1D_2$) and
\begin{equation}\label{expre of Btttpm tuta} \tilde{\tilde{\tilde{B}}}^{\pm}_{1}=\frac{1}{a^{\pm}_{1}}D_2\tilde{\tilde{B}}^{\pm}_2D_2,\qquad\tilde{\tilde{\tilde{B}}}^{\pm}_{2}=\frac{1}{a^{\pm}_{1}}D_2\tilde{\tilde{B}}^{\pm}_3D_2-\Big(\frac{1}{a^{\pm}_{1}}D_2\tilde{\tilde{B}}^{\pm}_2\Big)^2D_2.
\end{equation}
Then \eqref{asy expan of 1st 0} follows by successively combining \eqref{inverse of Mtuta3},\eqref{Relat of M3 and M3 tuta},\eqref{Relat of M2 tuta and M3},\eqref{Relat of M2 and M2 tuta},\eqref{Relat of M1 tuta and M2},\eqref{Relat of M1 and M1 tuta},\eqref{Relat of M tuta and M1} and \eqref{Relat of M and M tuta}.
\vskip0.2cm
{\bf{(iii) ${\bm{0}}$ is a second kind resonance of $H$ and {$\bm{\beta>27}$}.}} Take $N=9$ in \eqref{Puiseux expan of R0 0} and repeat {\bf{\underline{Step1}}} to {\bf{\underline{Step3}}} in {\bf{(i)}} and {\bf{(ii)}}. Since $T_1$ is not invertible on $S_1\ell^2(\Z)$, we need to further analyze the $\widetilde{M}^{\pm}_3(\mu)$ in {\bf{\underline{Step3}}} of {\bf{(ii)}}. Since ${\rm Ker}T_1\big|_{S_1\ell^2(\Z)}=S_2\ell^2(\Z)$, Lemma \ref{lemm of expa} yields the following further reduction.
\vskip0.15cm
{\bf{\underline{Step4:}}} \ $\widetilde{M}_{3}^{\pm}(\mu)$ is invertible on $S_1\ell^{2}(\Z)$ if and only if $$ M^{\pm}_{4}(\mu):=S_2-S_2\big(\widetilde{M}_{3}^{\pm}(\mu)+S_2\big)^{-1}S_2$$ is invertible on $S_2\ell^{2}(\Z)$. In this case,
\begin{equation}\label{Relat of M3 tuta and M4}
\big(\widetilde{M}_{3}^{\pm}(\mu)\big)^{-1}=\big(\widetilde{M}_{3}^{\pm}(\mu)+S_2\big)^{-1}\Big(I+S_2\big(M^{\pm}_{4}(\mu)\big)^{-1}S_2\big(\widetilde{M}_{3}^{\pm}(\mu)+S_2\big)^{-1}\Big).
\end{equation}
Let $D_2$ be the inverse of $T_1+S_2$ on $S_1\ell^2(\Z)$.
Then 
\begin{equation*}
\left(\widetilde{M}_{3}^{\pm}(\mu)+S_2\right)^{-1}=D_2-\sum\limits_{k=1}^{8}\mu^k\tilde{\tilde{\tilde{B}}}^{\pm}_{k}+D_2\Gamma_9(\mu)D_2, \quad \mu\rightarrow0.
\end{equation*}
We claim that  \begin{equation}\label{S2Bttt1S2,S2Bttt2S2}
S_2\tilde{\tilde{\tilde{B}}}^{\pm}_{1}S_2=0, \qquad S_2\tilde{\tilde{\tilde{B}}}^{\pm}_{2}S_2=d^{\pm}T_2,\qquad d^{\pm}=\pm8i.
\end{equation}
Assuming this relation for the moment, and using $S_2D_2=S_2=D_2S_2$, we obtain
\begin{equation}\label{rel betw M4 and wideM4}
M^{\pm}_4(\mu)=d^{\pm}\mu^2\Big(T_2+\frac{1}{d^{\pm}}\sum\limits_{k=3}^{8}\mu^{k-2}S_2\tilde{\tilde{\tilde{B}}}^{\pm}_{k}S_2+S_2\Gamma_7(\mu)S_2\Big)=:d^{\pm}\mu^2\widetilde{M}^{\pm}_4(\mu).
\end{equation}
By Remark \ref{remark of charac from inverti of operators}, the invertibility of $T_2$ on $S_2\ell^2(\Z)$ implies that $\widetilde{M}^{\pm}_4(\mu)$ is invertible on $S_2\ell^2(\Z)$, and
\begin{align}\label{inverse of Mtuta4}
\big(\widetilde{M}_{4}^{\pm}(\mu)\big)^{-1}
&=D_3+\sum\limits_{k=1}^{6}\mu^k\hat{B}^{\pm}_{k}+D_3\Gamma_7(\mu)D_3,\qquad D_3=T^{-1}_2,\qquad \mu\rightarrow0.
\end{align}
Thus, \eqref{asy expan 2nd 0} follows by successively combining \eqref{inverse of Mtuta4}, \eqref{rel betw M4 and wideM4}, \eqref{Relat of M3 tuta and M4}, \eqref{Relat of M3 and M3 tuta}, \eqref{Relat of M2 tuta and M3}, \eqref{Relat of M2 and M2 tuta}, \eqref{Relat of M1 tuta and M2}, \eqref{Relat of M1 and M1 tuta}, \eqref{Relat of M tuta and M1}, and \eqref{Relat of M and M tuta}.

Now, we verify \eqref{S2Bttt1S2,S2Bttt2S2}. The first identity follows from \eqref{expre of Btttpm tuta}, \eqref{expre of Bpm ttuta}, \eqref{expre of Bpm tuta}, \eqref{expr of Bkpm}, together with the following relations:
\begin{align*}
S_2D_j&=S_2=D_jS_2~(j=0,1,2),\quad D_0Q=D_0=QD_0,\quad  S_0D_1=D_1=D_1S_0,\\
S_2B^{\pm}_4S_2&=0,\quad S_2B^{\pm}_2Q=0=QB^{\pm}_2S_2,\quad S_2B^{\pm}_3S_0=0=S_0B^{\pm}_3S_2.
\end{align*}
For the second, it follows from \eqref{expre of Btttpm tuta} that 
\begin{align*}
S_2\tilde{\tilde{\tilde{B}}}^{\pm}_{2}S_2=\frac{1}{a^{\pm}_1}S_2\tilde{\tilde{B}}^{\pm}_3S_2-\Big(\frac{1}{a^{\pm}_{1}}\Big)^2S_2\tilde{\tilde{B}}^{\pm}_2D_2\tilde{\tilde{B}}^{\pm}_2S_2.
\end{align*}
Since $S_2\widetilde{B}^{\pm}_2D_1=0=D_1\widetilde{B}^{\pm}_2S_2$, $S_2B^{\pm}_2Q=0=QB^{\pm}_2S_2$, and $D_1D_2=D_2=D_2D_1$, we can compute
\begin{align}
\begin{split}
\frac{1}{a^{\pm}_1}S_2\tilde{\tilde{B}}^{\pm}_3S_2&=\frac{1}{a^{\pm}_1}S_2vG^{\pm}_3vS_2-\frac{1}{a^{\pm}_1a^{\pm}}S_2T^2_0S_2-\frac{1}{a^{\pm}_1a^{\pm}_{-1}}S_2vG^{\pm}_1vD_0vG^{\pm}_1vS_2,\\
\Big(\frac{1}{a^{\pm}_{1}}\Big)^2S_2\tilde{\tilde{B}}^{\pm}_2D_2\tilde{\tilde{B}}^{\pm}_2S_2&=\frac{1}{{\big(a^{\pm}_1\big)}^2}\Big(\frac{1}{a^{\pm}}S_2T_0vG^{\pm}_{-1}vD_0+\frac{1}{a^{\pm}_{-1}}S_2vG^{\pm}_1vD_0T_0D_0\Big)D_2\\
&\quad\times\Big(\frac{1}{a^{\pm}}D_0vG^{\pm}_{-1}vT_0S_2+\frac{1}{a^{\pm}_{-1}}D_0T_0D_0vG^{\pm}_1vS_2\Big),
\end{split}
\end{align}
which gives that $S_2\tilde{\tilde{\tilde{B}}}^{\pm}_{2}S_2=d^{\pm}T_2$ and this proves \eqref{S2Bttt1S2,S2Bttt2S2}.
\vskip0.2cm
{\bf{(iv) ${\bm{16}}$ is a regular point of $H$ and {$\bm{\beta>7}$}.}} Take $N=1$ in \eqref{Puiseux expan of R0 2}. Then as $\mu\rightarrow0$,
\begin{equation*}
M^{\pm}(2-\mu)=U+\tilde{v}JR^{\pm}_0((2-\mu)^4)J\tilde{v}
=\frac{d^{\pm}_{-1}}{\mu^{\frac{1}{2}}}\Big(\widetilde{P}+\frac{\mu^{\frac{1}{2}}}{d^{\pm}_{-1}}\widetilde{T}_0+\frac{\mu}{d^{\pm}_{-1}}\tilde{v}\widetilde{G}^{\pm}_{1}\tilde{v}+\Gamma_{\frac{3}{2}}(\mu)\Big)=:\frac{d^{\pm}_{-1}}{\mu^{\frac{1}{2}}}\widetilde{M}^{\pm}(\mu),
\end{equation*}
where we used the fact that $\tilde{v}\widetilde{G}^{\pm}_{-1}\tilde{v}=d^{\pm}_{-1}\widetilde{P}$ with $d^{\pm}_{-1}=\frac{\pm i}{32}\|V\|_{\ell^1}$. Then
\begin{align}\label{Relat of M and M tuta at 2} \left(M^{\pm}(2-\mu)\right)^{-1}=\frac{\mu^{\frac{1}{2}}}{d^{\pm}_{-1}}\big(\widetilde{M}^{\pm}(\mu)\big)^{-1}.
\end{align}

{\bf{\underline{Step1:}}} By Lemma \ref{lemm of expa}, $\widetilde{M}^{\pm}(\mu)$ is invertible on $\ell^{2}(\Z)$ if and only if $ M^{\pm}_1(\mu):=\widetilde{Q}-\widetilde{Q}\big(\widetilde{M}^{\pm}(\mu)+\widetilde{Q}\big)^{-1}\widetilde{Q}$ is invertible on $\widetilde{Q}\ell^{2}(\Z)$ and in this case, one has
\begin{align}\label{Relat of M tuta and M1 at 2}
\big(\widetilde{M}^{\pm}(\mu)\big)^{-1}=\big(\widetilde{M}^{\pm}(\mu)+\widetilde{Q}\big)^{-1}\Big[I+\widetilde{Q}\left(M^{\pm}_1(\mu)\right)^{-1}\widetilde{Q}\big(\widetilde{M}^{\pm}(\mu)+\widetilde{Q}\big)^{-1}\Big].
\end{align}
As $\mu\rightarrow0$, the Neumann series yields
\begin{equation}\label{expr of Ckpm}
\big(\widetilde{M}^{\pm}(\mu)+\widetilde{Q}\big)^{-1}=I-\sum\limits_{k=1}^{2}\mu^{\frac{k}{2}}C^{\pm}_{k}+\Gamma_{\frac{3}{2}}(\mu),\qquad C^{\pm}_{1}=\frac{1}{d^{\pm}_{-1}}\widetilde{T}_0,\quad C^{\pm}_{2}=\frac{1}{d^{\pm}_{-1}}\tilde{v}\widetilde{G}^{\pm}_1\tilde{v}-\Big(\frac{1}{d^{\pm}_{-1}}\widetilde{T}_0\Big)^2.
\end{equation}
Therefore,
\begin{equation}\label{Relat of M1 and M1 tuta at 2}
M^{\pm}_1(\mu)=\frac{\mu^{\frac{1}{2}}}{d^{\pm}_{-1}}\Big(\widetilde{Q}\widetilde{T}_0\widetilde{Q}+d^{\pm}_{-1}\mu^{\frac{1}{2}}\widetilde{Q}C^{\pm}_{2}\widetilde{Q}+\widetilde{Q}\Gamma_{1}(\mu)\widetilde{Q}\Big)=:\frac{\mu^{\frac{1}{2}}}{d^{\pm}_{-1}}\widetilde{M}^{\pm}_1(\mu).
\end{equation}
By Remark \ref{remark of charac from inverti of operators}, $\widetilde{Q}\widetilde{T}_0\widetilde{Q}$ is invertible on $\widetilde{Q}\ell^{2}(\Z)$. Denote its inverse by $E_0$. Then  $E_0\widetilde{Q}=E_0=\widetilde{Q}E_0$ and the Neumann series yields
\begin{equation}\label{inverse of M1 2}
\left( M^{\pm}_1(\mu)\right)^{-1}=\frac{d^{\pm}_{-1}}{\mu^{\frac{1}{2}}}E_0-\left(d^{\pm}_{-1}\right)^2E_0C^{\pm}_2E_0+E_0\Gamma_{\frac{1}{2}}(\mu)E_0,\quad\mu\rightarrow0.
\end{equation}
Then, combining \eqref{inverse of M1 2}, \eqref{Relat of M tuta and M1 at 2}, and \eqref{Relat of M and M tuta at 2}, we obtain the desired expansion \eqref{asy expan regular 2}.
\vskip0.1cm
{\bf{(v) ${\bm{16}}$ is a resonance of $H$ and {$\bm{\beta>11}$}.}} Take $N=3$ in \eqref{Puiseux expan of R0 2} and repeat {\bf{\underline{Step1}}} in {\bf{(iv)}}. Since ${\rm Ker}\widetilde{Q}\widetilde{T}_0\widetilde{Q}\big|_{\widetilde{Q}\ell^2(\Z)}=\widetilde{S}_0\ell^2(\Z)$, Lemma \ref{lemm of expa} yields the following further reduction.
\vskip0.15cm
{\bf{\underline{Step2:}}}\ $\widetilde{M}_{1}^{\pm}(\mu)$ is invertible on $\widetilde{Q}\ell^{2}(\Z)$ if and only if $ M^{\pm}_{2}(\mu):=\widetilde{S}_0-\widetilde{S}_0\big(\widetilde{M}_{1}^{\pm}(\mu)+\widetilde{S}_0\big)^{-1}\widetilde{S}_0$ is invertible on $\widetilde{S}_0\ell^{2}(\Z)$. In this case,
\begin{equation}\label{Relat of M1 tuta and M2 at 2}
\big(\widetilde{M}_{1}^{\pm}(\mu)\big)^{-1}=\big(\widetilde{M}_{1}^{\pm}(\mu)+\widetilde{S}_0\big)^{-1}\Big(I+\widetilde{S}_0\big(M^{\pm}_{2}(\mu)\big)^{-1}\widetilde{S}_0\big(\widetilde{M}_{1}^{\pm}(\mu)+\widetilde{S}_0\big)^{-1}\Big),
\end{equation}
where
\begin{align*}
\big(\widetilde{M}_{1}^{\pm}(\mu)+\widetilde{S}_0\big)^{-1}=E_0-\sum\limits_{k=1}^{3}\mu^{\frac{k}{2}}\widetilde{C}^{\pm}_{k}+E_0\Gamma_{2}(\mu)E_0,\qquad E_0=(\widetilde{Q}\widetilde{T}_0\widetilde{Q}+\widetilde{S}_0)^{-1},
\end{align*}
with 
\begin{align}\label{Ctutak}
\begin{split}
\widetilde{C}^{\pm}_{1}&=d^{\pm}_{-1}E_0C^{\pm}_{2}E_0,\quad \widetilde{C}^{\pm}_{2}=d^{\pm}_{-1}E_0C^{\pm}_{3}E_0-\big(d^{\pm}_{-1}\big)^2\big(E_0C^{\pm}_2\big)^2E_0,\\
C^{\pm}_{3}&=\frac{1}{d^{\pm}_{-1}}\tilde{v}\widetilde{G}^{\pm}_2\tilde{v}-\Big(\frac{1}{d^{\pm}_{-1}}\Big)^2\Big(\widetilde{T}_0\tilde{v}\widetilde{G}^{\pm}_1\tilde{v}+\tilde{v}\widetilde{G}^{\pm}_1\tilde{v}\widetilde{T}_0\Big)+\Big(\frac{1}{d^{\pm}_{-1}}\widetilde{T}_0\Big)^3.
\end{split}
\end{align}
Using $\widetilde{S}_0E_0=\widetilde{S}_0=E_0\widetilde{S}_0$, we obtain
\begin{equation}\label{Relat of M2 and M2 tuta at 2}
M^{\pm}_2(\mu)=d^{\pm}_1\mu^{\frac{1}{2}}\Big(\widetilde{T}_1+\frac{1}{d^{\pm}_{1}}\sum\limits_{k=2}^{3}\mu^{\frac{k-1}{2}}\widetilde{S}_0\widetilde{C}_{k}^{\pm}\widetilde{S}_0+\widetilde{S}_0\Gamma_{\frac{3}{2}}(\mu)\widetilde{S}_0\Big)=:d^{\pm}_1\mu^{\frac{1}{2}}\widetilde{M}^{\pm}_2(\mu),\quad d^{\pm}_1=\pm i.
\end{equation}
By Remark \ref{remark of charac from inverti of operators}, $\widetilde{T}_1$ is invertible on $\widetilde{S}_0\ell^2(\Z)$. Using Von-Neumann expansion yields
\begin{equation}\label{inverse of Mtuta2}
\big(\widetilde{M}_{2}^{\pm}(\mu)\big)^{-1}=E_1-\mu^{\frac{1}{2}}\tilde{\tilde{C}}^{\pm}_{1}-\mu \tilde{\tilde{C}}^{\pm}_2+E_1\Gamma_{\frac{3}{2}}(\mu)E_1,\qquad \tilde{\tilde{C}}^{\pm}_{1}=\frac{1}{d^{\pm}_1}E_1\widetilde{C}^{\pm}_2E_1,\quad \mu\rightarrow0,
\end{equation}
where $E_1$ is the inverse of $\widetilde{T}_1$ on $\widetilde{S}_0\ell^2(\Z)$.
Combining \eqref{inverse of Mtuta2}, \eqref{Relat of M2 and M2 tuta at 2}, \eqref{Relat of M1 tuta and M2 at 2}, \eqref{Relat of M1 and M1 tuta at 2}, \eqref{Relat of M tuta and M1 at 2} and \eqref{Relat of M and M tuta at 2}, \eqref{asy expan resonance 2} is established. 
\vskip0.2cm
{\bf{(vi) ${\bm{16}}$ is an eigenvalue of $H$ and {$\bm{\beta>15}$}.}} Take $N=5$ in \eqref{Puiseux expan of R0 2} and repeat {\bf{\underline{Step1}}} in {\bf{(iv)}} and {\bf{\underline{Step2}}} in {\bf{(v)}}. Since ${\rm Ker}\widetilde{T}_1\big|_{\widetilde{S}_0\ell^2(\Z)}=\widetilde{S}_1\ell^2(\Z)$, Lemma \ref{lemm of expa} yields the following further reduction.
\vskip0.15cm
{\bf{\underline{Step3:}}}\ $\widetilde{M}_{2}^{\pm}(\mu)$ is invertible on $\widetilde{S}_0\ell^{2}(\Z)\Leftrightarrow M^{\pm}_{3}(\mu):=\widetilde{S}_1-\widetilde{S}_1\big(\widetilde{M}_{2}^{\pm}(\mu)+\widetilde{S}_1\big)^{-1}\widetilde{S}_1$ is invertible on $\widetilde{S}_1\ell^{2}(\Z)$. In this case,
\begin{equation}\label{Relat of M2 tuta and M3 at 2}
\big(\widetilde{M}_{2}^{\pm}(\mu)\big)^{-1}=\big(\widetilde{M}_{2}^{\pm}(\mu)+\widetilde{S}_1\big)^{-1}\Big(I+\widetilde{S}_1\big(M^{\pm}_{3}(\mu)\big)^{-1}\widetilde{S}_1\big(\widetilde{M}_{2}^{\pm}(\mu)+\widetilde{S}_1\big)^{-1}\Big),
\end{equation}
where
\begin{align*}
\big(\widetilde{M}_{2}^{\pm}(\mu)+\widetilde{S}_1\big)^{-1}=E_1-\sum\limits_{k=1}^{4}\mu^{\frac{k}{2}}\tilde{\tilde{C}}^{\pm}_{k}+E_1\Gamma_{5/2}(\mu)E_1,\quad \mu\rightarrow0,
\end{align*}
with $E_1=(\widetilde{T}_1+\widetilde{S}_1)^{-1}$ on $\widetilde{S}_0\ell^2(\Z)$ and $\tilde{\tilde{C}}^{\pm}_{1}$ defined in \eqref{inverse of Mtuta2}. Using $\widetilde{T}_0\widetilde{S}_1=0=\widetilde{Q}C^{\pm}_2\widetilde{S}_1$ and $E_1\widetilde{S}_1=\widetilde{S}_1=\widetilde{S}_1E_1$, we have
\begin{equation*}
M^{\pm}_3(\mu)=\frac{1}{d^{\pm}_1}\mu^{\frac{1}{2}}\Big(\widetilde{T}_2+{d^{\pm}_1}\sum\limits_{k=2}^{4}\mu^{\frac{k-1}{2}}\widetilde{S}_1\tilde{\tilde{C}}^{\pm}_{k}\widetilde{S}_1+\widetilde{S}_1\Gamma_{2}(\mu)\widetilde{S}_1\Big)=:\frac{\mu^{\frac{1}{2}}}{d^{\pm}_1}\widetilde{M}^{\pm}_3(\mu).
\end{equation*}
Since $\widetilde{T}_2$ is invertible on $\widetilde{S}_1\ell^2(\Z)$, it follows that $\widetilde{M}^{\pm}_3(\mu)$ is invertible on $\widetilde{S}_1\ell^2(\Z)$ for sufficiently small $\mu>0$. Applying the Neumann series and successively combining \eqref{Relat of M2 tuta and M3 at 2}, \eqref{Relat of M2 and M2 tuta at 2}, \eqref{Relat of M1 tuta and M2 at 2}, \eqref{Relat of M1 and M1 tuta at 2}, \eqref{Relat of M tuta and M1 at 2}, and \eqref{Relat of M and M tuta at 2}, we obtain the desired expansion \eqref{asy expan eigenvalue 2}. This completes the proof.
\end{proof}

\begin{thebibliography}{99}
\bibitem[32]{JN01} A.~Jensen and G.~Nenciu, \newblock A unified approach to resolvent expansions at thresholds, \newblock {\em Rev. Math. Phys.}~{\bf13}~(2001), no. 6, 717--754. 
\bibitem[34]{KKK06} A.~I. Komech, E.~A. Kopylova and M.~Kunze, \newblock Dispersive estimates for 1{D} discrete {S}chr\"odinger and {K}lein-{G}ordon equations, \newblock {\em Appl. Anal.}~{\bf85}~(2006), no. 12, 1487--1508. 
\end{thebibliography}
\end{document}
